\chapter{What the historical evidence actually establishes}\label{ch:evidence}
A research programme becomes more credible when it preserves the conditions
under which it failed or produced a flat result. Replacing every older
chapter with a success-oriented summary would lose that information. This
edition therefore retains selected original evidence chapters as historical
source pages and adds a careful reading guide.

\section{Read a result together with its model}
Appendix H-III reproduces original Chapters 53--59, covering D1--D10,
Gate 1B, adversarial/service-alignment work and O14. The old potential,
regeneration, preservation rules and selectors remain attached to those
accounts. Their numerical results were not recomputed under the new formula.

The appropriate question is not ``Did EBU pass?'' without qualification.
Ask which equation, physical domain, comparison, horizon and falsifier were
tested. A test of quote-code conformance and a test of service alignment
answer different questions even when both involve the same numerical quote.

\section{Why the old distinctions still teach useful lessons}
The original soft-versus-hard preservation comparisons distinguish a penalty
from a permission boundary. That lesson survives a change in potential:
assigning a large burden does not physically prevent an action. The
observational-quote comparison distinguishes logging from influencing the
plant. The adversarial examples ask whether a locally closing account can
still coexist with unrepresented or persistent harm.

These are methodological lessons. They are not evidence that the Gaussian
capacity policy has passed analogous tests. A new actor, menu or state
representation changes the tested system and requires its own registered
evaluation.

\section{O14 and a later finalized gate are different records}
The original O14 chapter preserves its own capability and flat-outcome
analysis. It should not be relabelled as Gate 1D-C or as Gaussian evidence.
Some older analytical documents also retain operational status text from
before the later gate was completed. Their mathematical content can be
useful without treating their historical status paragraphs as current.

The later committed Gate 1D-C manifest records a finalized study
\cite{gate}. Its execution SHA is \texttt{8793bda}; its inventory contains
thirty registered runs and six thousand trace rows. The recorded stdout
specifies three worlds, two timestep settings and five arms, with two
hundred ticks per run. Its reported outcome counts are twenty
\texttt{physical\_impossibility}, eight
\texttt{safe\_rationing\_physical\_scarcity}, and two
\texttt{preserve\_and\_serve}; it reports no fired falsifier.

Those are narrow historical source facts, not thirty independent proofs
of the accounting theory or a long-run Gaussian result. The revised books
do not rerun, tune, extend or reinterpret the frozen campaign. The source
manifest and original scientific protocol govern its actual claims.

\section{A comparison can be invalid even when its numbers are real}
An arm allowed several simultaneous deliveries has a different capability
set from an arm restricted to one action per configured source. Such an
arm can be informative as a capability benchmark, but it is not automatically
a matched test of the contribution of the EBU selector. Historical
comparison restrictions must remain visible when reading an outcome table.

Likewise, a new Gaussian arm with a different physical menu cannot claim
an advantage of capacity accounting alone unless that difference is part
of the registered estimand. The artifact may be genuine while the proposed
cross-model inference is unsupported.

\section{What preservation of the appendix means}
The appendix is reproduced source material, not a fresh independent audit
of every archived run. Its figures and statements keep their original
scope and publication-era wording. The new source manifest identifies the
exact PDF pages. Readers can return to the full preserved originals for
cross-references outside the extract.

This treatment avoids two opposite errors: discarding useful history
because the active model changed, and making the new model appear validated
by repainting that history. Neither is necessary for a coherent series.

\section*{Chapter recap}
Historical experiments remain evidence about their own registered systems.
They guide questions and expose pitfalls for new research, but they cannot
answer an unexecuted Gaussian comparison by analogy alone.
